% recovery-paper.tex
% Recovering Write-Protected NVMe SSDs via USB Bridge XRAM Command Injection
% J. Sullivan, 2026

\documentclass[10pt,twocolumn]{article}

% ── Packages ──────────────────────────────────────────────────────────
\usepackage[utf8]{inputenc}
\usepackage[T1]{fontenc}
\usepackage{amsmath,amssymb}
\usepackage{graphicx}
\usepackage{listings}
\usepackage{hyperref}
\usepackage{booktabs}
\usepackage{xcolor}
\usepackage[margin=0.75in]{geometry}
\usepackage{float}
\usepackage{caption}
\usepackage{subcaption}
\usepackage{enumitem}
\usepackage{tikz}
\usepackage{tabularx}
\usepackage{fancyhdr}

% ── Colors ────────────────────────────────────────────────────────────
\definecolor{codebg}{HTML}{F5F5F5}
\definecolor{codeframe}{HTML}{CCCCCC}
\definecolor{codecomment}{HTML}{6A737D}
\definecolor{codestring}{HTML}{032F62}
\definecolor{codekeyword}{HTML}{D73A49}
\definecolor{hexblue}{HTML}{0366D6}
\definecolor{successgreen}{HTML}{28A745}
\definecolor{failred}{HTML}{CB2431}

% ── Listings Configuration ────────────────────────────────────────────
\lstset{
  basicstyle=\ttfamily\footnotesize,
  backgroundcolor=\color{codebg},
  frame=single,
  rulecolor=\color{codeframe},
  breaklines=true,
  breakatwhitespace=false,
  columns=fullflexible,
  keepspaces=true,
  tabsize=2,
  showstringspaces=false,
  commentstyle=\color{codecomment}\itshape,
  keywordstyle=\color{codekeyword}\bfseries,
  stringstyle=\color{codestring},
  numbers=left,
  numberstyle=\tiny\color{gray},
  numbersep=5pt,
  xleftmargin=1em,
  framexleftmargin=1em,
  captionpos=b,
}

\lstdefinestyle{hex}{
  basicstyle=\ttfamily\footnotesize\color{hexblue},
  numbers=none,
  xleftmargin=0.5em,
  framexleftmargin=0.5em,
}

\lstdefinestyle{shell}{
  language=bash,
  morekeywords={sudo,zig,modprobe,rmmod,echo,dd,xxd,wipefs},
}

% ── Hyperref Configuration ────────────────────────────────────────────
\hypersetup{
  colorlinks=true,
  linkcolor=black,
  citecolor=hexblue,
  urlcolor=hexblue,
  pdftitle={Recovering Write-Protected NVMe SSDs via USB Bridge XRAM Command Injection},
  pdfauthor={J. Sullivan},
}

% ── Headers ───────────────────────────────────────────────────────────
\pagestyle{fancy}
\fancyhf{}
\fancyhead[L]{\footnotesize Sullivan -- NVMe Recovery via XRAM Injection}
\fancyhead[R]{\footnotesize\thepage}
\renewcommand{\headrulewidth}{0.4pt}

% ── Title ─────────────────────────────────────────────────────────────
\title{%
  \textbf{Recovering Write-Protected NVMe SSDs via\\
  USB Bridge XRAM Command Injection}
}
\author{
  J.~Sullivan\\
  \texttt{github.com/jsullivan2/hiberpower-ntfs}
}
\date{March 2026}

% ══════════════════════════════════════════════════════════════════════
\begin{document}
\maketitle

% ── Abstract ──────────────────────────────────────────────────────────
\begin{abstract}
We present a novel technique for recovering NVMe SSDs that have entered
firmware-level write-protection mode following Flash Translation Layer (FTL)
corruption. When an NVMe SSD is connected through a USB bridge, the bridge
firmware's opcode whitelist typically blocks all recovery commands (Format NVM,
Sanitize, Security Send/Receive), making the drive appear unrecoverable via
USB. We demonstrate that the ASMedia ASM2362 USB-to-NVMe bridge chip exposes
vendor-specific SCSI commands (0xE4/0xE5) providing direct read/write access
to its internal XRAM, where the NVMe Admin Submission Queue resides. By
writing NVMe command structures directly into this queue memory and ringing
the PCIe doorbell via TLP generation, we bypass the bridge's opcode whitelist
entirely. Using this method, we successfully recovered a Phison PS5012-E12
controller (Silicon Power 256GB) that had been in read-only mode for over a
month, executing a Sanitize Block Erase command that standard tools could not
deliver. The technique generalizes to other ASMedia bridge variants and
potentially to other bridge chipsets with exposed internal memory.
\end{abstract}

\tableofcontents

% ══════════════════════════════════════════════════════════════════════
\section{Introduction}
\label{sec:intro}

NVMe solid-state drives can enter a firmware-controlled read-only mode when
internal errors corrupt the Flash Translation Layer (FTL). This protective
mechanism preserves existing data when the controller detects that further
write operations could lead to data loss. While the NVMe specification
provides recovery commands---Format NVM (opcode \texttt{0x80}) and Sanitize
(opcode \texttt{0x84})---these commands require direct admin queue access that
is not always available.

A common scenario occurs when the affected SSD is housed in a USB enclosure.
USB-to-NVMe bridge chips translate between the SCSI command set (used over
USB) and NVMe admin commands. However, these bridges implement firmware-level
whitelists that restrict which NVMe opcodes can pass through the translation
layer. The ASMedia ASM2362---one of the most widely deployed USB 3.1 Gen 2
NVMe bridges---only forwards \texttt{Identify} (opcode \texttt{0x06}) and
\texttt{Get Log Page} (opcode \texttt{0x02}) through its proprietary 0xE6
passthrough CDB. All other NVMe admin commands, including those necessary for
drive recovery, are silently dropped.

This paper describes a technique we developed and successfully used to
recover a Silicon Power 256GB NVMe SSD (Phison PS5012-E12 controller) that
had been in write-protected mode since January 2026 following a Windows
hibernate corruption event combined with an unexpected power loss. The core
insight is that the ASM2362 exposes vendor-specific SCSI commands
(\texttt{0xE4} for XRAM read, \texttt{0xE5} for XRAM write) that provide
direct access to the bridge chip's internal memory, where the NVMe Admin
Submission Queue (Admin SQ) resides. By writing a 64-byte NVMe Submission
Queue Entry (SQE) directly into the Admin SQ via XRAM and then ringing the
SQ Tail Doorbell via PCIe Transaction Layer Packet (TLP) generation, we can
inject arbitrary NVMe admin commands that bypass the bridge firmware's opcode
filter entirely.

\subsection{Contributions}

This work makes the following contributions:

\begin{enumerate}[leftmargin=*]
  \item \textbf{XRAM injection technique}: A method for injecting arbitrary
    NVMe admin commands through a USB bridge by writing directly to the
    bridge's internal XRAM, bypassing the firmware opcode whitelist.
  \item \textbf{PCIe TLP doorbell mechanism}: Documentation of how to ring
    the NVMe Admin SQ Tail Doorbell via the ASM2362's PCIe TLP generation
    engine without disconnecting USB.
  \item \textbf{Practical recovery procedure}: A step-by-step procedure that
    successfully recovered a write-protected Phison PS5012-E12 SSD through
    USB, including the discovery that the Admin SQ depth is 4 (not 8) and
    that only Sanitize Block Erase (not Format NVM or Crypto Erase) clears
    the FTL write protection on this controller.
  \item \textbf{Open-source implementation}: A Zig-based tool
    (\texttt{asm2362-tool}) comprising approximately 4,000 lines of code
    with 35 unit tests implementing the complete injection pipeline.
\end{enumerate}

% ══════════════════════════════════════════════════════════════════════
\section{Background}
\label{sec:background}

\subsection{NVMe Architecture}
\label{sec:nvme-arch}

The NVMe (Non-Volatile Memory Express) specification~\cite{nvme-spec-2.0}
defines a register-level interface between the host and NVMe controller. At
its core, NVMe uses a paired-queue model: Submission Queues (SQ) for sending
commands from host to controller, and Completion Queues (CQ) for receiving
results. Queue 0 is the Admin Queue pair, used for device management; queues
1--N are I/O Queue pairs for data transfer.

Each Submission Queue Entry (SQE) is 64 bytes, organized as 16 DWORDs
(32-bit words, little-endian). Table~\ref{tab:sqe-format} shows the SQE
structure. The host writes command entries to the SQ and notifies the
controller by writing the new tail pointer to the SQ Tail Doorbell register,
located at offset \texttt{0x1000} from BAR0 for the Admin SQ.

\begin{table}[htbp]
\centering
\caption{NVMe Submission Queue Entry (64 bytes)}
\label{tab:sqe-format}
\footnotesize
\begin{tabular}{@{}llp{4.0cm}@{}}
\toprule
\textbf{Offset} & \textbf{Field} & \textbf{Description} \\
\midrule
\texttt{00--03} & CDW0  & Opcode [7:0], FUSE [9:8], PSDT [15:14], CID [31:16] \\
\texttt{04--07} & NSID  & Namespace Identifier \\
\texttt{08--0F} & Rsvd  & Reserved (8 bytes) \\
\texttt{10--13} & MPTR  & Metadata Pointer (low 32) \\
\texttt{14--17} & MPTR  & Metadata Pointer (high 32) \\
\texttt{18--1B} & PRP1  & PRP Entry 1 (low 32) \\
\texttt{1C--1F} & PRP1  & PRP Entry 1 (high 32) \\
\texttt{20--23} & PRP2  & PRP Entry 2 (low 32) \\
\texttt{24--27} & PRP2  & PRP Entry 2 (high 32) \\
\texttt{28--2B} & CDW10 & Command-specific DWORD 10 \\
\texttt{2C--2F} & CDW11 & Command-specific DWORD 11 \\
\texttt{30--33} & CDW12 & Command-specific DWORD 12 \\
\texttt{34--37} & CDW13 & Command-specific DWORD 13 \\
\texttt{38--3B} & CDW14 & Command-specific DWORD 14 \\
\texttt{3C--3F} & CDW15 & Command-specific DWORD 15 \\
\bottomrule
\end{tabular}
\end{table}

The Completion Queue Entry (CQE) is 16 bytes and includes a Status Field that
indicates command success or failure, along with the Command ID (CID) for
correlating completions with submissions.

\subsection{USB-NVMe Bridges}
\label{sec:usb-bridges}

USB-to-NVMe bridge chips translate between the SCSI command set used by the
USB Mass Storage and UAS (USB Attached SCSI) protocols and the NVMe command
set used by the SSD. The ASMedia ASM2362 is among the most common such
bridges, deployed in numerous USB 3.1 Gen 2 (10 Gbps) NVMe enclosures.

\subsubsection{ASM2362 Architecture}

The ASM2362 is built around an 8051-compatible microcontroller core running
at approximately 114.3 MHz~\cite{cyrozap-usb-pcie}. Key architectural
features include:

\begin{itemize}[leftmargin=*,nosep]
  \item \textbf{XRAM}: 64KB of mapped internal memory, serving as both
    working memory for the 8051 firmware and as DMA-mapped regions for
    NVMe queue structures.
  \item \textbf{PCIe interface}: Gen 3 x2 lanes connecting to the NVMe
    device, with a TLP generation engine accessible from XRAM registers.
  \item \textbf{SPI flash}: External flash storage for firmware, with no
    signature verification on firmware updates.
  \item \textbf{UART debug}: 921600 baud 8N1 at 3.3V on pins 62/63.
\end{itemize}

\subsubsection{The 0xE6 Passthrough CDB}
\label{sec:e6-passthrough}

The ASM2362 implements a vendor-specific 16-byte SCSI CDB with opcode
\texttt{0xE6} for tunneling NVMe admin commands over USB. This is the
mechanism used by smartmontools~\cite{smartmontools} (via the
\texttt{sntasmedia\_device} class) and Silicon Power's SP Toolbox
application. Table~\ref{tab:e6-cdb} shows the CDB byte layout.

\begin{table}[htbp]
\centering
\caption{ASM2362 0xE6 Passthrough CDB Format (16 bytes)}
\label{tab:e6-cdb}
\footnotesize
\begin{tabular}{@{}clp{3.8cm}@{}}
\toprule
\textbf{Byte} & \textbf{Field} & \textbf{Description} \\
\midrule
0     & \texttt{0xE6}      & ASMedia passthrough opcode \\
1     & NVMe OPC           & NVMe admin opcode \\
2     & Reserved           & \texttt{0x00} \\
3     & CDW10[7:0]         & Low byte of CDW10 \\
4--5  & Reserved           & \texttt{0x00 0x00} \\
6     & CDW10[23:16]       & Bits 23--16 of CDW10 \\
7     & CDW10[31:24]       & High byte of CDW10 \\
8--11 & CDW13              & Big-endian \\
12--15& CDW12              & Big-endian \\
\bottomrule
\end{tabular}
\end{table}

A critical limitation of this CDB format is that CDW10 bits [15:8] and
CDW11, CDW14, CDW15 cannot be transmitted. More importantly, the bridge
firmware implements an \textbf{opcode whitelist} that only forwards two NVMe
admin opcodes:

\begin{itemize}[leftmargin=*,nosep]
  \item \texttt{0x02} -- Get Log Page
  \item \texttt{0x06} -- Identify
\end{itemize}

All other opcodes---including Format NVM (\texttt{0x80}), Sanitize
(\texttt{0x84}), Set Features (\texttt{0x09}), Get Features (\texttt{0x0A}),
Security Send (\texttt{0x81}), and Security Receive (\texttt{0x82})---are
silently dropped at the bridge firmware level. The SCSI layer returns
\texttt{GOOD} status, but the command never reaches the NVMe controller.

\subsection{Phison PS5012-E12 Controller}
\label{sec:phison}

The Phison PS5012-E12 is a quad-core, eight-channel NVMe 1.3 controller
widely used in consumer SSDs including Silicon Power, Kingston, Corsair, and
Sabrent products. Key characteristics of the specific drive under
investigation:

\begin{itemize}[leftmargin=*,nosep]
  \item VID: \texttt{0x1E4B} (Phison Electronics)
  \item Firmware: \texttt{H211011a}
  \item OACS: \texttt{0x0006} (Format NVM + Firmware Download supported)
  \item SANICAP: \texttt{0x00000002} (Block Erase only)
  \item Advertised capacity: 256 GB (500,118,192 sectors at 512B)
\end{itemize}

\subsection{Related Work}
\label{sec:related}

The foundational reverse engineering of ASMedia USB bridge chips was
performed by cyrozap~\cite{cyrozap-usb-pcie}, who documented the XRAM memory
map, vendor SCSI commands (0xE0--0xE8), and PCIe TLP generation mechanism.
The smartmontools project~\cite{smartmontools} implements the 0xE6
passthrough for SMART monitoring. The ASMTool project~\cite{asmtool}
provides firmware dumping capabilities. Phison-specific recovery tools exist
through the usbdev.ru community~\cite{phison-reinitial}, but require Windows
and vendor-specific firmware matching. Professional recovery via PC-3000
SSD~\cite{acelab-phison} is available but expensive (\$300--\$1,500).

Prior work on SSD power fault resilience~\cite{zheng2013understanding} has
shown that 13 out of 15 consumer SSDs exhibit data loss during power fault
testing, establishing the prevalence of the failure mode that initiated our
investigation.

% ══════════════════════════════════════════════════════════════════════
\section{Problem Analysis}
\label{sec:problem}

\subsection{Failure Scenario}

The target drive experienced a combined Windows hibernate and unexpected power
loss event in January 2026. After the event, the drive exhibited the
following behavior:

\begin{itemize}[leftmargin=*,nosep]
  \item SCSI READ commands succeed (data is readable)
  \item SCSI WRITE commands report success but data is unchanged
  \item NVMe admin commands via 0xE6 return ``Medium not present''
    (Sense Key \texttt{0x02}, ASC \texttt{0x3A})
  \item No operating system write-protect flags are set
    (\texttt{blockdev --getro} returns 0)
\end{itemize}

\subsection{Silent Write Failure Characterization}

The drive's silent write failure was confirmed through three independent
tests, all of which demonstrated that write operations appear to succeed at
the SCSI protocol level but produce no change on the media:

\begin{lstlisting}[style=shell,caption={Test 1: dd to sector 0}]
$ sudo dd if=/dev/zero of=/dev/sdb \
    bs=512 count=1 conv=fsync oflag=direct
512 bytes copied, 0.000313659 s, 1.6 MB/s

$ sudo xxd -l 16 /dev/sdb
00000000: 33c0 8ed0 bc00 7c8e ...
# MBR boot code remains -- NOT zeroed
\end{lstlisting}

\begin{lstlisting}[style=shell,caption={Test 2: wipefs signature erasure}]
$ sudo wipefs --all --force /dev/sdb
/dev/sdb: 8 bytes were erased at offset 0x200

$ sudo xxd -s 512 -l 8 /dev/sdb
00000200: 4546 4920 5041 5254
# "EFI PART" header still present
\end{lstlisting}

Standard Linux tools---\texttt{fdisk}, \texttt{gdisk},
\texttt{blkdiscard}, \texttt{sg\_format}---all reported success but left the
drive data unchanged. The SCSI Write Protect mode page correctly reported
\texttt{WP=0}, and Linux's \texttt{blockdev --getro} returned 0,
confirming that no host-side write protection was active.

\subsection{SMART Analysis}

The SMART/Health log, retrieved through the 0xE6 passthrough (opcode
\texttt{0x02}, which is whitelisted), revealed:

\begin{table}[htbp]
\centering
\caption{SMART/Health log data (pre-recovery)}
\label{tab:smart-pre}
\footnotesize
\begin{tabular}{@{}lr@{}}
\toprule
\textbf{Field} & \textbf{Value} \\
\midrule
Critical Warning            & \texttt{0x00} \\
Available Spare             & 1\% \\
Available Spare Threshold   & 10\% \\
Media and Data Integrity Errors & 292 \\
Unsafe Shutdowns            & 270 \\
Power On Hours              & 5,936 \\
Data Units Written          & 31,289,178 \\
\bottomrule
\end{tabular}
\end{table}

A key finding was that \textbf{SMART Critical Warning was \texttt{0x00}}---bit
3 (``media placed in read-only mode'') was \textbf{not} set. This contradicted
our initial hypothesis that the SMART bit 3 flag was driving the write
protection. The controller had placed the media in read-only mode through an
internal mechanism not reflected in the standard SMART critical warning byte.

\subsection{The 0xE6 Whitelist Discovery}

Through systematic testing of all NVMe admin opcodes observed in SP Toolbox
USB captures, we confirmed the bridge firmware whitelist. Table~\ref{tab:whitelist}
shows the results.

\begin{table}[htbp]
\centering
\caption{ASM2362 0xE6 opcode whitelist testing}
\label{tab:whitelist}
\footnotesize
\begin{tabular}{@{}clcc@{}}
\toprule
\textbf{Opcode} & \textbf{Command} & \textbf{Passes?} & \textbf{Evidence} \\
\midrule
\texttt{0x02} & Get Log Page     & \color{successgreen}Yes & Data returned \\
\texttt{0x06} & Identify         & \color{successgreen}Yes & Data returned \\
\texttt{0x09} & Set Features     & \color{failred}No  & Silent drop \\
\texttt{0x0A} & Get Features     & \color{failred}No  & Silent drop \\
\texttt{0x80} & Format NVM       & \color{failred}No  & Silent drop \\
\texttt{0x81} & Security Send    & \color{failred}No  & Silent drop \\
\texttt{0x82} & Security Receive & \color{failred}No  & Silent drop \\
\texttt{0x84} & Sanitize         & \color{failred}No  & Silent drop \\
\bottomrule
\end{tabular}
\end{table}

The ``silent drop'' behavior is particularly insidious: the SCSI layer
returns \texttt{GOOD} status for blocked commands, making it appear as though
the NVMe controller accepted and executed the command. Only by observing the
unchanged drive state could we determine that the commands never reached the
controller.

% ══════════════════════════════════════════════════════════════════════
\section{XRAM Injection Methodology}
\label{sec:methodology}

\subsection{ASM2362 Vendor SCSI Commands}

Beyond the 0xE6 NVMe passthrough, the ASM2362 exposes additional
vendor-specific SCSI commands~\cite{cyrozap-usb-pcie}. Table~\ref{tab:vendor-cmds}
lists the complete set.

\begin{table}[htbp]
\centering
\caption{ASM2362 vendor-specific SCSI commands}
\label{tab:vendor-cmds}
\footnotesize
\begin{tabular}{@{}cllc@{}}
\toprule
\textbf{Opcode} & \textbf{Command} & \textbf{Direction} & \textbf{Size} \\
\midrule
\texttt{0xE0} & Read Config  & Dev$\to$Host & 128B \\
\texttt{0xE1} & Write Config & Host$\to$Dev & 128B \\
\texttt{0xE2} & Flash Read   & Dev$\to$Host & Variable \\
\texttt{0xE3} & Firmware Write & Host$\to$Dev & Variable \\
\texttt{0xE4} & XDATA Read   & Dev$\to$Host & 1--255B \\
\texttt{0xE5} & XDATA Write  & None         & 1B (in CDB) \\
\texttt{0xE6} & NVMe Admin   & Dev$\to$Host & Variable \\
\texttt{0xE8} & Reset        & None         & N/A \\
\bottomrule
\end{tabular}
\end{table}

The 0xE4 and 0xE5 commands provide direct read/write access to the
bridge chip's 64KB XRAM address space. Critically, the NVMe Admin Submission
Queue is DMA-mapped into XRAM, meaning that XRAM writes are reflected in the
queue structures that the NVMe controller reads.

\subsection{UAS vs.\ BOT Mode Requirement}
\label{sec:uas-bot}

An essential prerequisite discovered during implementation is that the vendor
SCSI commands \texttt{0xE4}, \texttt{0xE5}, and \texttt{0xE8}
\textbf{require the \texttt{usb-storage} driver (BOT mode)}, not the
\texttt{uas} driver. Under UAS, these commands return \texttt{DID\_ERROR}
with \texttt{host\_status=0x07} (errno \texttt{-75 EOVERFLOW}).

To switch from UAS to BOT mode:

\begin{lstlisting}[style=shell,caption={Switching ASM2362 to BOT mode}]
sudo bash -c '
  echo "1-3" > /sys/bus/usb/drivers/usb/unbind
  rmmod uas
  rmmod usb_storage
  modprobe usb-storage quirks=174c:2362:u
  echo "1-3" > /sys/bus/usb/drivers/usb/bind
'
\end{lstlisting}

The device re-enumerates (typically as \texttt{/dev/sdc}). Verification:
\begin{lstlisting}[style=shell,caption={Verifying BOT mode}]
readlink /sys/devices/.../1-3:1.0/driver
# Should show: usb-storage (not uas)
\end{lstlisting}

\subsection{XRAM Memory Map}
\label{sec:xram-map}

Table~\ref{tab:xram-map} shows the XRAM memory regions relevant to NVMe
command injection, as documented by cyrozap~\cite{cyrozap-usb-pcie} and
refined through our probing.

\begin{table}[htbp]
\centering
\caption{ASM2362 XRAM memory map (relevant regions)}
\label{tab:xram-map}
\footnotesize
\begin{tabular}{@{}llp{3.2cm}@{}}
\toprule
\textbf{Address} & \textbf{Size} & \textbf{Contents} \\
\midrule
\texttt{0x07F0}      & 6B    & Firmware version \\
\texttt{0xA000--AFFF} & 4KB   & NVMe I/O Submission Queue \\
\texttt{0xB000--B0FF} & 256B  & Admin SQ (4 entries$\times$64B) \\
\texttt{0xB100--B1FF} & 256B  & Admin SQ (unused region) \\
\texttt{0xB200--B296} & 151B  & PCIe MMIO / TLP engine \\
\texttt{0xB800--B9FF} & 512B  & I/O Completion Queue \\
\texttt{0xBC00--BC7F} & 128B  & Admin CQ (8$\times$16B) \\
\texttt{0xF000--FFFF} & 4KB   & NVMe data buffer \\
\bottomrule
\end{tabular}
\end{table}

A significant discovery was that the \textbf{Admin SQ depth is 4, not 8}.
Although the XRAM region \texttt{0xB000--0xB1FF} provides 512 bytes (space
for 8 entries), only slots 0--3 are within the queue boundary as configured
by the bridge firmware. This was confirmed by observing SQ Head (SQHD)
wrapping behavior in Completion Queue entries: the SQHD value wraps at 4,
not 8. Commands injected into slots 4--7 were never processed by the
controller.

Similarly, the \textbf{Admin CQ base is at \texttt{0xBC00}}, not
\texttt{0xB800} as initially assumed from the documentation. The region at
\texttt{0xB800} contains I/O Completion Queue entries (SQID=1).

\subsection{CDB Formats for XRAM Access}
\label{sec:cdb-formats}

The vendor SCSI CDB formats, corrected from our implementation experience
(the original documentation had byte positions transposed):

\vspace{0.5em}
\noindent\textbf{0xE4 -- XDATA Read (6-byte CDB):}
\begin{lstlisting}[style=hex]
Byte 0: E4           ; Opcode
Byte 1: <length>     ; Bytes to read (1-255)
Byte 2: 00           ; Padding
Byte 3: <addr_hi>    ; Address bits [15:8]
Byte 4: <addr_lo>    ; Address bits [7:0]
Byte 5: 00           ; Padding
Direction: FROM_DEV, transfer = <length> bytes
\end{lstlisting}

\noindent\textbf{0xE5 -- XDATA Write (6-byte CDB):}
\begin{lstlisting}[style=hex]
Byte 0: E5           ; Opcode
Byte 1: <value>      ; Single byte to write
Byte 2: 00           ; Padding
Byte 3: <addr_hi>    ; Address bits [15:8]
Byte 4: <addr_lo>    ; Address bits [7:0]
Byte 5: 00           ; Padding
Direction: NONE (byte is embedded in the CDB itself)
\end{lstlisting}

\noindent\textbf{0xE8 -- Reset (12-byte CDB):}
\begin{lstlisting}[style=hex]
Byte 0:  E8          ; Opcode
Byte 1:  <type>      ; 0x00=CPU reset, 0x01=PCIe reset
Bytes 2-11: 00       ; Padding
Direction: NONE
\end{lstlisting}

Note that the 0xE5 write command embeds the data byte \emph{within the CDB
itself}---there is no SCSI data transfer phase. This means writing a 64-byte
NVMe SQE requires issuing 64 separate 0xE5 SCSI commands, each writing one
byte. While slow, this is reliable and provides byte-level write verification
capability via 0xE4 readback.

\subsection{PCIe TLP Doorbell Mechanism}
\label{sec:pcie-tlp}

After writing the NVMe command into the Admin SQ, the SQ Tail Doorbell must
be rung to notify the NVMe controller of the new submission. The ASM2362's
PCIe TLP generation engine allows crafting arbitrary PCIe Transaction Layer
Packets from XRAM registers. The procedure, ported from
cyrozap's \texttt{pcie\_gen\_req()} function~\cite{cyrozap-usb-pcie}:

\begin{enumerate}[leftmargin=*,nosep]
  \item Read BAR0 from PCIe config space via Type 1 Configuration Read TLP
    (bus=1, dev=0, fn=0, offset=\texttt{0x10}). For our device:
    BAR0 = \texttt{0x00D00000}.
  \item Admin SQ Tail Doorbell address = BAR0 + \texttt{0x1000} =
    \texttt{0x00D01000}.
  \item Write the new tail value (big-endian) to the TLP data register at
    XRAM \texttt{0xB220}.
  \item Write the 12-byte TLP header (3$\times$ u32, big-endian) to XRAM
    \texttt{0xB210}:
    \begin{itemize}[nosep]
      \item DW0: \texttt{0x40000001} (Memory Write, 1 DW payload)
      \item DW1: Byte enable mask
      \item DW2: Masked doorbell address
    \end{itemize}
  \item Clear the timeout bit: write \texttt{0x01} to XRAM \texttt{0xB296}.
  \item Trigger the TLP: write \texttt{0x0F} to XRAM \texttt{0xB254}.
  \item Poll XRAM \texttt{0xB296} bit 2 until ready.
  \item Send the TLP: write \texttt{0x04} to XRAM \texttt{0xB296}.
\end{enumerate}

This is a \emph{posted} memory write TLP, meaning no completion is expected.
The USB connection stays alive throughout---unlike using the 0xE8 PCIe reset
command, which disconnects USB.

An important limitation: \textbf{PCIe memory \emph{read} TLPs do not work}
through this engine (they hang/timeout). Only config reads and posted memory
writes function correctly. This means we cannot read the NVMe controller's
registers (e.g., CSTS, CQ entries) via PCIe TLP; we can only read what is
mapped into XRAM.

% ══════════════════════════════════════════════════════════════════════
\section{Failed Approaches}
\label{sec:failures}

Before arriving at the successful injection technique, numerous approaches
were attempted and failed. Table~\ref{tab:failures} catalogs these dead
ends.

\begin{table*}[htbp]
\centering
\caption{Summary of failed recovery approaches}
\label{tab:failures}
\footnotesize
\begin{tabularx}{\textwidth}{@{}lXX@{}}
\toprule
\textbf{Approach} & \textbf{Method} & \textbf{Failure Reason} \\
\midrule
0xE6 Format NVM &
  Build 0xE6 CDB with opcode \texttt{0x80} &
  Bridge firmware whitelist silently drops the command; SCSI returns GOOD but NVMe controller never receives it \\
\addlinespace
0xE6 Sanitize &
  Build 0xE6 CDB with opcode \texttt{0x84} &
  Same whitelist: silently dropped \\
\addlinespace
0xE6 Set Features &
  Attempt to clear write protection via FID \texttt{0x84} (Namespace Write Protect) &
  Whitelist blocks opcode \texttt{0x09}; also, NWPC is a separate mechanism from FTL read-only \\
\addlinespace
0xE6 Security Send/Recv &
  ATA password protocol (\texttt{0xEF}) to clear security state &
  Whitelist blocks opcodes \texttt{0x81}/\texttt{0x82} \\
\addlinespace
XRAM inject to slots 4--7 &
  Write NVMe command to SQ entries 4--7 &
  Admin SQ depth = 4 (not 8); entries beyond slot 3 are outside the queue boundary and never consumed \\
\addlinespace
Sanitize Crypto Erase &
  Inject Sanitize with SANACT=4 (Crypto Erase) &
  SANICAP = \texttt{0x00000002}: only Block Erase is supported; Crypto Erase not available on this controller \\
\addlinespace
Format NVM via inject &
  Inject Format NVM (\texttt{0x80}) to slot 0 with doorbell &
  Controller accepts command (CQ shows success), but FTL write protection persists; Format alone does not clear the protection state \\
\addlinespace
PCIe reset as doorbell &
  Use 0xE8 type=1 (PCIe reset) to trigger SQ processing &
  Disconnects USB; also, Format NVM executed during reset caused 0B capacity (reverted on power cycle) \\
\addlinespace
Wine + SP Toolbox &
  Run Silicon Power's Windows tool under Wine on Linux &
  Wine does not implement \texttt{IOCTL\_SCSI\_PASS\_THROUGH\_DIRECT} (\texttt{0x4D014}); all device operations return \texttt{STATUS\_NOT\_SUPPORTED} \\
\addlinespace
SMART bit 3 hypothesis &
  Assume Critical Warning bit 3 drives read-only mode and can be cleared &
  SMART Critical Warning = \texttt{0x00}; bit 3 was not set; the original hypothesis was wrong \\
\bottomrule
\end{tabularx}
\end{table*}

\subsection{The Format NVM Paradox}

An instructive failure was the Format NVM command via XRAM injection. The
command was accepted by the controller---the Completion Queue showed success
status (SC=0, SCT=0) with the correct Command ID---yet the drive remained in
read-only mode afterward. This reveals that the Phison PS5012-E12's FTL
write protection is not cleared by a Format NVM operation, even one that
succeeds at the NVMe protocol level. Only the Sanitize command, which
performs a more thorough NAND block erase including over-provisioned areas,
was able to reset the FTL state.

\subsection{The Queue Depth Discovery}

Injecting commands into slots 4--7 produced no Completion Queue entries and
no observable effect. By dumping the Admin CQ (\texttt{0xBC00}) after
injecting to various slots, we observed that the SQHD (SQ Head) field in
completion entries never exceeded 3, confirming a depth of 4. This
contradicts the naive assumption based on the 512-byte XRAM allocation
(which would suggest 8 entries of 64 bytes each).

% ══════════════════════════════════════════════════════════════════════
\section{Successful Recovery}
\label{sec:recovery}

\subsection{SANICAP Analysis}

Before attempting the Sanitize command, we verified the controller's
Sanitize Capabilities through the Identify Controller data (retrieved via
the whitelisted 0xE6 Identify command):

\begin{lstlisting}[style=hex,caption={Identify Controller -- Sanitize Capabilities}]
OACS    = 0x0006  ; Format NVM + Firmware Download
SANICAP = 0x00000002
  Bit 0 (Crypto Erase):    0  ; Not supported
  Bit 1 (Block Erase):     1  ; SUPPORTED
  Bit 2 (Overwrite):       0  ; Not supported
\end{lstlisting}

Only Sanitize Block Erase (SANACT=2) was available. Crypto Erase
(SANACT=4), which we had initially attempted, was not supported by this
controller.

\subsection{Injection Procedure}

The successful recovery used the following exact command sequence:

\subsubsection{Step 1: Switch to BOT Mode}

\begin{lstlisting}[style=shell]
sudo modprobe usb-storage quirks=174c:2362:u
# Verify: /dev/sda via usb-storage driver
\end{lstlisting}

\subsubsection{Step 2: Probe and Verify Bridge}

\begin{lstlisting}[style=shell]
sudo ./zig-out/bin/asm2362-tool probe /dev/sda
# Confirms ASM2362 detected, firmware version
\end{lstlisting}

\subsubsection{Step 3: Dry-Run Injection}

\begin{lstlisting}[style=shell]
sudo ./zig-out/bin/asm2362-tool inject \
  --inject-cmd=sanitize-block \
  --slot=0 --tail=1 --cid=0x4242 \
  --dry-run /dev/sda
\end{lstlisting}

The dry-run writes the SQE to XRAM and verifies it via readback, but
\textbf{does not ring the doorbell}. This allows visual inspection of the
injected command before committing.

\subsubsection{Step 4: Live Injection}

\begin{lstlisting}[style=shell]
sudo ./zig-out/bin/asm2362-tool inject \
  --inject-cmd=sanitize-block \
  --slot=0 --tail=1 --cid=0x4242 \
  --force /dev/sda
\end{lstlisting}

This writes the 64-byte NVMe SQE to Admin SQ slot 0 at XRAM
\texttt{0xB000}, verifies via readback, then rings the Admin SQ Tail
Doorbell at BAR0+\texttt{0x1000} = \texttt{0x00D01000} with tail value 1
via a posted PCIe memory write TLP.

The injected SQE (Sanitize Block Erase, CID=\texttt{0x4242}):

\begin{lstlisting}[style=hex,caption={Sanitize Block Erase SQE (64 bytes, LE)}]
B000: 84 00 42 42 00 00 00 00  ; CDW0: OPC=0x84 CID=0x4242
B008: 00 00 00 00 00 00 00 00  ; Reserved, MPTR
B010: 00 00 00 00 00 00 00 00  ; MPTR(hi), PRP1
B018: 00 00 00 00 00 00 00 00  ; PRP1(hi), PRP2
B020: 00 00 00 00 00 00 00 00  ; PRP2(hi)
B028: 02 00 00 00 00 00 00 00  ; CDW10=0x02 (SANACT=2)
B030: 00 00 00 00 00 00 00 00  ; CDW12-CDW13
B038: 00 00 00 00 00 00 00 00  ; CDW14-CDW15
\end{lstlisting}

\subsubsection{Step 5: Verify Completion}

\begin{lstlisting}[style=shell]
sudo ./zig-out/bin/asm2362-tool admin-cq /dev/sda
\end{lstlisting}

The Admin CQ at XRAM \texttt{0xBC00} showed a new completion entry:

\begin{lstlisting}[style=hex,caption={Admin CQ entry confirming success}]
BC00: xx xx xx xx 42 42 xx xx  ; CID=0x4242
BC08: xx xx xx xx 00 00 01 xx  ; Status=0x0000, Phase=1
\end{lstlisting}

Key fields: CID=\texttt{0x4242} (matches our injection), Status Code
Type (SCT) = 0, Status Code (SC) = 0 (success), Phase Tag = 1 (valid
new entry).

\subsubsection{Step 6: Bridge Reset and Re-enumeration}

\begin{lstlisting}[style=shell]
sudo ./zig-out/bin/asm2362-tool reset \
  --reset-type=0 /dev/sda
\end{lstlisting}

A CPU reset (type=0) forces the 8051 to restart, which re-initializes the
NVMe controller from the now-erased NAND. The device re-enumerates on the
USB bus.

\subsubsection{Step 7: Sanitize Status Verification}

After re-enumeration, the Sanitize Log (Get Log Page with LID=\texttt{0x81})
confirmed:

\begin{table}[htbp]
\centering
\caption{Post-sanitize verification}
\label{tab:post-sanitize}
\footnotesize
\begin{tabular}{@{}lr@{}}
\toprule
\textbf{Field} & \textbf{Value} \\
\midrule
SSTAT (Sanitize Status) & \texttt{0x0001} (Complete) \\
Write test (1MB random data) & Success, 17.7 MB/s \\
Read-back verification & MD5 match confirmed \\
Identify Controller & Responds normally \\
\bottomrule
\end{tabular}
\end{table}

\subsection{Post-Recovery Write Verification}

\begin{lstlisting}[style=shell,caption={Write verification after recovery}]
# Write 1MB of random data
$ sudo dd if=/dev/urandom of=/dev/sda \
    bs=1M count=1 oflag=direct
1+0 records out, 1048576 bytes, 17.7 MB/s

# Read back and verify
$ sudo dd if=/dev/sda bs=1M count=1 | md5sum
# Hash matches -- writes are functional
\end{lstlisting}

The drive was fully operational after recovery: reads, writes, and NVMe
admin commands all functioned normally. All previous data was erased by
the Sanitize Block Erase operation.

% ══════════════════════════════════════════════════════════════════════
\section{Implementation}
\label{sec:implementation}

\subsection{Tool Architecture}

The recovery tool (\texttt{asm2362-tool}) is implemented in
Zig~\cite{zig-lang} (version 0.13.0+) and comprises approximately 4,000
lines of code with 35 unit tests. The architecture follows a layered
design:

\begin{lstlisting}[style=hex,caption={Source tree structure}]
src/
  main.zig            CLI entry point
  log.zig             JSON command logging
  scsi/
    sg_io.zig         Linux SG_IO ioctl wrapper
    sense.zig         SCSI sense data parsing
  asm2362/
    passthrough.zig   0xE6 CDB (Identify, SMART)
    commands.zig      NVMe command helpers
    xram.zig          0xE4/0xE5/0xE8 + injection
  nvme/
    identify.zig      Identify Controller/NS
  frida/
    hooks.js          Windows DeviceIoControl capture
    replay.zig        Frida capture replay
  analysis/
    probe.zig         Capability detection
\end{lstlisting}

\subsection{SG\_IO Interface}

All SCSI commands are issued via the Linux \texttt{SG\_IO}
ioctl~\cite{linux-sg-io}, which provides synchronous SCSI command passthrough
to block devices. The \texttt{sg\_io.zig} module wraps the
\texttt{sg\_io\_hdr} structure:

\begin{lstlisting}[language=C,caption={Key sg\_io\_hdr fields used}]
struct sg_io_hdr {
  int interface_id;     // 'S' for SCSI
  int dxfer_direction;  // SG_DXFER_FROM_DEV, etc.
  unsigned char *cmdp;  // CDB pointer
  unsigned char cmd_len;// CDB length (6, 12, 16)
  void *dxferp;         // Data buffer pointer
  unsigned int dxfer_len;// Data buffer length
  unsigned char *sbp;   // Sense buffer pointer
  unsigned char mx_sb_len;// Max sense length
  unsigned int timeout; // Milliseconds
  // ... additional fields
};
\end{lstlisting}

\subsection{Key Data Structures}

The \texttt{NvmeSqEntry} structure models the 64-byte NVMe SQE with
serialization methods:

\begin{lstlisting}[language=C,caption={NvmeSqEntry (simplified)}]
pub const NvmeSqEntry = struct {
    cdw0: u32,  // OPC | CID<<16
    nsid: u32,
    reserved8: u32 = 0,
    reserved12: u32 = 0,
    mptr_lo: u32 = 0,
    mptr_hi: u32 = 0,
    prp1_lo: u32 = 0,
    prp1_hi: u32 = 0,
    prp2_lo: u32 = 0,
    prp2_hi: u32 = 0,
    cdw10: u32 = 0,
    cdw11: u32 = 0,
    cdw12: u32 = 0,
    cdw13: u32 = 0,
    cdw14: u32 = 0,
    cdw15: u32 = 0,

    pub fn toBytes(self) [64]u8;
    pub fn fromBytes(bytes: []const u8) NvmeSqEntry;
};
\end{lstlisting}

Command crafting functions build SQEs for specific NVMe commands:
\texttt{craftFormatNvmEntry()}, \texttt{craftSanitizeEntry()}, and
\texttt{craftSetFeaturesEntry()}.

\subsection{Injection Pipeline}

The \texttt{injectCommand()} function implements the six-phase injection
pipeline:

\begin{enumerate}[leftmargin=*,nosep]
  \item \textbf{Read} current Admin SQ state via 0xE4 (512 bytes from
    \texttt{0xB000}).
  \item \textbf{Select} SQ slot: explicit (via \texttt{--slot}) or
    auto-detect first empty entry.
  \item \textbf{Write} 64 bytes via 64 individual 0xE5 commands, each
    verified via 0xE4 readback.
  \item \textbf{Verify} full 64-byte readback against expected SQE bytes.
  \item \textbf{Doorbell}: Ring Admin SQ Tail Doorbell via PCIe TLP
    (skipped in dry-run mode).
  \item \textbf{Post-read} Admin SQ and CQ state for status verification.
\end{enumerate}

\subsection{Safety Mechanisms}

The tool implements multiple safety layers:

\begin{itemize}[leftmargin=*,nosep]
  \item \textbf{Dry-run default}: All injection commands default to dry-run
    mode; \texttt{--force} is required for live execution.
  \item \textbf{Write verification}: Every XRAM byte write is verified via
    readback before proceeding to the next byte.
  \item \textbf{Address validation}: Write operations are restricted to
    Admin SQ (\texttt{0xB000--B1FF}) and data buffer (\texttt{0xF000--FFFF})
    regions.
  \item \textbf{CID correlation}: User-specified Command IDs
    (\texttt{--cid}) enable unambiguous CQ entry correlation.
  \item \textbf{Explicit slot/tail}: Queue slot and doorbell tail can be
    specified explicitly to avoid assumptions about queue state.
\end{itemize}

% ══════════════════════════════════════════════════════════════════════
\section{Discussion}
\label{sec:discussion}

\subsection{Implications for SSD Data Recovery}

The XRAM injection technique demonstrates that USB-to-NVMe bridges, despite
their firmware-level command restrictions, can be bypassed through direct
memory manipulation. This has significant implications for the data recovery
industry, where USB-connected drives are common in consumer scenarios:

\begin{itemize}[leftmargin=*,nosep]
  \item Drives previously considered unrecoverable via USB may be salvageable
    without physical disassembly.
  \item The technique is non-destructive to the bridge hardware---only XRAM
    is modified, and this is volatile (lost on reset).
  \item Professional recovery services using PC-3000
    SSD~\cite{acelab-phison,rossmann-ssd} (\$300--\$1,500) may not be
    necessary for certain failure modes.
\end{itemize}

\subsection{Generalizability}

The technique should generalize to other ASMedia bridge variants (ASM2362,
ASM2364, ASM236X family) that share the same 8051 architecture and vendor
SCSI command set. The XRAM memory map may differ, but the discovery
methodology (systematic 0xE4 probing) remains applicable.

Other bridge chipsets (JMicron JMS583, Realtek RTL9210) use different
architectures and may or may not expose equivalent internal memory access.
This remains an area for future investigation.

\subsection{Ethical Considerations}

The XRAM injection technique bypasses security controls in the USB bridge
firmware. While our use case---recovering a personally owned drive---is
benign, the same technique could potentially be misused to:

\begin{itemize}[leftmargin=*,nosep]
  \item Bypass security erase verification (making forensic sanitization
    unreliable through USB bridges).
  \item Inject malicious NVMe commands into drives connected through shared
    USB infrastructure.
  \item Exploit bridge firmware vulnerabilities for firmware-level attacks.
\end{itemize}

We note that the ASM2362 already lacks firmware signature
verification~\cite{cyrozap-usb-pcie}, making these concerns largely moot
for this specific chip---the security boundary was never robust. However,
responsible disclosure of bridge-level bypasses remains important as bridge
firmware security matures.

\subsection{Limitations}

\begin{enumerate}[leftmargin=*,nosep]
  \item \textbf{BOT mode required}: The vendor SCSI commands fail under UAS,
    requiring a driver switch that may not be possible in all environments.
  \item \textbf{Bridge-specific XRAM layout}: Queue base addresses, depths,
    and MMIO register locations are specific to the ASM2362 firmware version
    and must be discovered per-device.
  \item \textbf{Single-byte writes}: The 0xE5 command writes one byte per
    SCSI transaction, making 64-byte SQE injection slow (64 SCSI commands
    plus 64 verification reads).
  \item \textbf{No PCIe memory reads}: The TLP engine does not support
    non-posted reads, preventing direct access to NVMe controller registers.
  \item \textbf{Data destruction}: Sanitize Block Erase destroys all data on
    the drive. This technique recovers the \emph{drive}, not the
    \emph{data}.
  \item \textbf{Firmware-specific behavior}: The Phison PS5012-E12's
    response to Sanitize vs.\ Format commands may differ from other
    controllers; the finding that only Sanitize clears the FTL protection
    may not generalize.
\end{enumerate}

% ══════════════════════════════════════════════════════════════════════
\section{Conclusion}
\label{sec:conclusion}

We have presented a novel technique for recovering write-protected NVMe SSDs
through USB bridge XRAM command injection. By exploiting the ASMedia ASM2362's
vendor-specific SCSI commands for direct XRAM access, we bypassed the bridge
firmware's NVMe opcode whitelist and successfully delivered a Sanitize Block
Erase command to a Phison PS5012-E12 controller that had been in read-only
mode for over a month.

The key technical insights are: (1)~the NVMe Admin Submission Queue resides
in bridge XRAM and can be written via vendor SCSI commands; (2)~the PCIe TLP
generation engine can ring the SQ Tail Doorbell without disconnecting USB;
(3)~vendor SCSI commands require BOT mode, not UAS; (4)~the Admin SQ depth
and CQ base address must be empirically determined, as they differ from
documentation; and (5)~Sanitize Block Erase, not Format NVM, is required to
clear FTL write protection on the PS5012-E12.

The open-source implementation (\texttt{asm2362-tool}) and this
documentation provide a reusable framework for NVMe SSD recovery through
ASMedia USB bridges, potentially saving drives that would otherwise require
expensive professional recovery services.

% ── Acknowledgments ──────────────────────────────────────────────────
\section*{Acknowledgments}

The author thanks cyrozap for the foundational ASM236x reverse engineering
work~\cite{cyrozap-usb-pcie} without which this technique would not have
been possible, the smartmontools developers~\cite{smartmontools} for
the reference \texttt{sntasmedia} implementation, and the usbdev.ru
community for documenting Phison controller recovery procedures.

% ── References ────────────────────────────────────────────────────────
\bibliographystyle{plain}
\bibliography{references}

% ══════════════════════════════════════════════════════════════════════
\appendix

\section{Hardware Stack Diagram}
\label{app:stack}

\begin{lstlisting}[style=hex,caption={Hardware stack (host to media)}]
+-------------------------------------------+
| Linux Host (Rocky Linux 10)               |
|   SG_IO ioctl -> /dev/sda                 |
+--------------------+----------------------+
                     | USB 3.1 Gen 2
+--------------------+----------------------+
| ASMedia ASM2362 USB-NVMe Bridge           |
|   VID:PID = 174c:2362                     |
|   CPU: 8051 @ 114.3 MHz                   |
|   XRAM: 64KB (Admin SQ at 0xB000)        |
|   Firmware whitelist: 0x02, 0x06 only     |
|                                           |
| BYPASS PATH:                              |
|   0xE4 -> Read XRAM                       |
|   0xE5 -> Write XRAM (Admin SQ)           |
|   0xB210-B296 -> PCIe TLP engine          |
|     -> Doorbell at BAR0+0x1000            |
+--------------------+----------------------+
                     | PCIe Gen 3 x2
+--------------------+----------------------+
| NVMe SSD: Phison PS5012-E12              |
|   VID: 0x1E4B                             |
|   FW: H211011a                            |
|   OACS: 0x0006                            |
|   SANICAP: 0x00000002 (Block Erase only)  |
|   State: FTL write-protected              |
+--------------------+----------------------+
                     | 8-channel NAND
+--------------------+----------------------+
| Silicon Power SPCC M.2 PCIe SSD          |
|   Capacity: 256 GB                        |
|   Media: 3D TLC NAND                      |
|   State: Read-only (FTL corrupted)        |
+-------------------------------------------+
\end{lstlisting}

\section{Complete Injection Workflow}
\label{app:workflow}

Figure~\ref{fig:workflow} shows the complete injection workflow as a
sequential procedure.

\begin{figure*}[htbp]
\centering
\begin{lstlisting}[style=hex,caption={Injection workflow},label=fig:workflow]
[1. Switch to BOT mode]
    |
    v
[2. Probe bridge (0xE4 read at 0x0000)]
    |
    v
[3. Read Admin SQ (0xE4, 256B at 0xB000)]
    |-- Parse 4 x 64-byte SQ entries
    |-- Identify target slot (0-3)
    v
[4. Craft NVMe SQE]
    |-- Sanitize: OPC=0x84, CDW10=0x02
    |-- Set CID=0x4242 for correlation
    v
[5. Write SQE to XRAM (64 x 0xE5 commands)]
    |-- Each byte: write + verify readback
    |-- Abort on any verification failure
    v
[6. Full readback verification]
    |-- 0xE4 read 64 bytes at slot address
    |-- Compare against expected SQE bytes
    v
[7. DRY-RUN GATE]
    |-- --dry-run: STOP HERE (safe inspection)
    |-- --force: CONTINUE to doorbell
    v
[8. Read BAR0 via PCIe config read TLP]
    |-- Type 1 cfg read, bus=1 dev=0 fn=0
    |-- BAR0 = 0x00D00000
    v
[9. Ring doorbell via PCIe memory write TLP]
    |-- Write tail=1 to BAR0+0x1000
    |-- Posted write: no completion expected
    |-- USB stays alive
    v
[10. Read Admin CQ (0xE4 at 0xBC00)]
    |-- Find entry with CID=0x4242
    |-- Check Status: SCT=0, SC=0 = success
    v
[11. CPU reset (0xE8 type=0)]
    |-- Bridge restarts, NVMe re-inits
    |-- Device re-enumerates on USB
    v
[12. Post-recovery verification]
    |-- Write 1MB random data
    |-- Read back and compare hashes
    |-- SMART log: SSTAT=0x0001
\end{lstlisting}
\end{figure*}

\section{NVMe Admin Opcodes Reference}
\label{app:opcodes}

\begin{table}[htbp]
\centering
\caption{NVMe Admin opcodes relevant to recovery}
\label{tab:opcodes-ref}
\footnotesize
\begin{tabular}{@{}clc@{}}
\toprule
\textbf{Opcode} & \textbf{Command} & \textbf{CDW10 Key Bits} \\
\midrule
\texttt{0x00} & Delete I/O SQ & QID \\
\texttt{0x01} & Create I/O SQ & QID, QSIZE \\
\texttt{0x02} & Get Log Page  & LID, NUMD \\
\texttt{0x04} & Delete I/O CQ & QID \\
\texttt{0x05} & Create I/O CQ & QID, QSIZE \\
\texttt{0x06} & Identify      & CNS \\
\texttt{0x09} & Set Features  & FID, SV \\
\texttt{0x0A} & Get Features  & FID, SEL \\
\texttt{0x80} & Format NVM    & LBAF, SES \\
\texttt{0x81} & Security Send & SECP, SPSP \\
\texttt{0x82} & Security Recv & SECP, SPSP \\
\texttt{0x84} & Sanitize      & SANACT \\
\bottomrule
\end{tabular}
\end{table}

\section{XRAM Hex Dump: Admin SQ Pre-Recovery}
\label{app:hexdump}

The Admin SQ state before injection, showing 4 occupied entries from the
bridge's initialization sequence (Set Features $\times$2, Create I/O CQ,
Create I/O SQ):

\begin{lstlisting}[style=hex,caption={Admin SQ at 0xB000 (pre-injection)}]
B000: 09 00 01 00 00 00 00 00  ; [0] Set Features
B008: 00 00 00 00 00 00 00 00  ;     CID=0x0001
B010: 00 00 00 00 00 00 00 00
B018: 00 00 00 00 00 00 00 00
B020: 00 00 00 00 00 00 00 00
B028: 07 00 00 00 FE 00 FE 00  ;     FID=0x07 (Num Queues)
B030: 00 00 00 00 00 00 00 00
B038: 00 00 00 00 00 00 00 00

B040: 09 00 02 00 00 00 00 00  ; [1] Set Features
B048: 00 00 00 00 00 00 00 00  ;     CID=0x0002
...
B080: 05 00 03 00 00 00 00 00  ; [2] Create I/O CQ
B088: 00 00 00 00 00 00 00 00  ;     CID=0x0003
...
B0C0: 01 00 04 00 00 00 00 00  ; [3] Create I/O SQ
B0C8: 00 00 00 00 00 00 00 00  ;     CID=0x0004
...
B100: 00 00 00 00 00 00 00 00  ; [4] (empty -- beyond
B108: 00 00 00 00 00 00 00 00  ;      queue boundary)
\end{lstlisting}

The four entries represent the bridge firmware's NVMe initialization
sequence: configure queue counts, create I/O completion queue, create I/O
submission queue. All four slots are occupied, which is why injection to
\texttt{slot=0} with \texttt{tail=1} was necessary---the slot is
overwritten with our Sanitize command, and the tail pointer wraps.

\end{document}
